% ===================================================================
%  TABLICA Nr 8 — Żniwa. — Polowanie. — Winobranie. — Rybołówstwo.
%  Meme regime que les precedents.
%
%  LE BLOC t08-c2-03-1 EST L'UN DES TROIS ECARTS DECLARES, et pour une
%  raison de COMPOSITION et non de sens : le fac-simile ido y ecrit
%  « 13) » sans sa parenthese ouvrante, si bien que le releve des
%  renvois ne le voit pas. Les traductions ecrivent le renvoi entier —
%  le neerlandais et le tcheque le font — et cette colonne aussi.
%
%  ET LA NOTE (*) DU TABLEAU PARLE DE L'IDO LUI-MEME : elle propose
%  « mesono » pour distinguer la moisson des autres recoltes, et dit
%  que l'Academie ne l'a pas adopte. On la traduit telle quelle, mot
%  pour mot : c'est un document sur la langue, pas une phrase du
%  livret.
% ===================================================================

%%K t08-noto-1 noto
\VUnotes{143.2mm}{%
\textit{(*) Ponieważ to słowo nie jest ścisłe: czy chodzi o zbiór lnu, winobranie, zbiór
jabłek? Może dobrze byłoby użyć «} \VUgras{mesono} \textit{» zaproponowanego w Mondo XIV,
s. 242. Ale dotąd ten nowy rdzeń nie został przyjęty przez Akademię i czeka na dyskusje
zgodnie ze zwyczajowym regulaminem idystycznym.}%
}

%%K t08-apar-1 apar
\VUsaut{5.0mm}
\VUtitre{15.0pt}{60}{TABLICA Nr 8}
\VUsaut{3.0mm}
\VUcentre{13.2pt}{90}{\textit{Scena pierwsza.}}
\VUsaut{3.0mm}
\VUpk{10.2pt}{40}{Żniwa (*).}
\VUsaut{4.0mm}
\VUpk{10.2pt}{40}{Widoki wsi.}
\VUsaut{3.0mm}

%%K t08-c1-01-1 p
1. --- Wraz z \VUgras{latem} raz jeszcze zmienił się widok wsi. Ziarno powierzone
bruździe wykiełkowało; zielona roślinka wyszła z ziemi i wypuściła kłos. Ziarno się
tworzyło, potem dojrzewało, a odtąd trzeba je już tylko zebrać.

%%K t08-c1-02-1 p
2. --- \VUgras{Kwiaty polne} są rozwinięte. \VUgras{Chabry} \textsuperscript{(1)},
\VUgras{maki} \textsuperscript{(2)} i \VUgras{naparstnice} \textsuperscript{(3)} mieszają
z jednostajnym odcieniem pól zbożowych wesoły kontrast swoich świeżych \VUgras{koron}.

%%K t08-c1-03-1 p
3. --- Drzewa owocowe okryły się liśćmi; owoc dojrzał. Mała \VUgras{czereśnia}
\textsuperscript{(4)} obfituje w piękne \VUgras{czereśnie} \textsuperscript{(5)}, które
dzieci z wielką przyjemnością zbierają i jedzą.

%%K t08-c1-04-1 p
4. --- Odtąd bydło może spędzać cały dzień na wolnym powietrzu.

%%K t08-c1-04-2 p
\VUgras{Jagnięta} \textsuperscript{(6)} urosły i stały się pięknymi \VUgras{owcami}
\textsuperscript{(7)} i pięknymi \VUgras{baranami} \textsuperscript{(8)} o gęstym wełnianym
runie. Spokojnie skubią \VUgras{trawę} \VUgras{pastwiska} \textsuperscript{(9)} upstrzonego
\VUgras{stokrotkami}.

%%K t08-c1-04-3 p
Wkrótce zacznie się strzyc \VUgras{te} zwierzęta, żeby mieć ich \VUgras{wełnę}, z której
wyrabia się sukno na nasze ubrania.

%%K t08-tit-2 sub
\VUsaut{5.0mm}
\VUpk{10.2pt}{40}{Pasterz.}
\VUsaut{3.4mm}

%%K t08-c1-05-1 p
5. --- \VUgras{Pasterz} \textsuperscript{(10)} zdaje się mało na nie zważać. Dba tylko o
burzę, która przechodzi na widnokręgu, a \VUgras{juhas} \textsuperscript{(13)}, który
przykłada \VUgras{lirę korbową} \textsuperscript{(14)} do warg, wydaje się jeszcze
mniej zatroskany.

%%K t08-c1-06-1 p
6. --- Upał jest przytłaczający; bo \VUgras{pies} \textsuperscript{(15)} też przychodzi
ugasić pragnienie; chłepcze świeżą wodę \VUgras{strumienia} \textsuperscript{(16)} i płoszy
biedną \VUgras{żabkę} \textsuperscript{(17)}, która przycupnęła w nadbrzeżnych trawach.
Ta skacze w czystą wodę, gdzie kwitną \VUgras{grzybienie} \textsuperscript{(18)}.

%%K t08-c1-07-1 p
7. --- \VUgras{Pliszka} \textsuperscript{(19)} kąpie się przy \VUgras{źródełku}
\textsuperscript{(20)}, które tryska między dwiema skałkami, i chowa się pod
\VUgras{ciernistymi krzewami} \textsuperscript{(21)} i \VUgras{krzakami głogu}
\textsuperscript{(22)}.

%%K t08-tit-3 sub
\VUsaut{5.0mm}
\VUpk{10.2pt}{40}{Żniwa.}
\VUsaut{3.4mm}

%%K t08-c1-08-1 p
8. --- Dla wiejskich dzieci żniwa są bardzo zabawną porą. Radośnie
\VUgras{koziołkują}\textsuperscript{(24)} na \VUgras{stertach słomy}
\textsuperscript{(23)}, pomagają \VUgras{żniwiarzom} \textsuperscript{(26)} i podczas gdy
ci koszą \VUgras{pole zbożowe} \textsuperscript{(27)} swoimi \VUgras{kosami}
\textsuperscript{(25)} dobrze naostrzonymi \VUgras{osełką} \textsuperscript{(30)}, one
zbierają zboże i wiążą je w \VUgras{snopy} \textsuperscript{(31)} albo w
\VUgras{wiązki} \textsuperscript{(32)}.

%%K t08-c1-09-1 p
9. --- Czasem pozwala się największym z nich ciąć \VUgras{sierpem}
\textsuperscript{(12)} \VUgras{złote źdźbła} \textsuperscript{(28)}, które niosą ciężkie
\VUgras{kłosy} \textsuperscript{(29)} pełne ziarna; a najmniejsze, pilnując
\VUgras{bydła} \textsuperscript{(33)}, które pasie się na \VUgras{pastwisku}
\textsuperscript{(34)} w cieniu wielkiej \VUgras{lipy} \textsuperscript{(35)}, chwytają
swoją \VUgras{siatką} \textsuperscript{(36)} piękne \VUgras{motyle}.

%%K t08-c1-09-2 p
Niektóre nawet wspinają się na \VUgras{wóz} \textsuperscript{(37)}, żeby układać snopy,
które podaje się im \VUgras{widłami} \textsuperscript{(38)}.

%%K t08-c1-10-1 p
10. --- Na całej \VUgras{równinie} \textsuperscript{(39)}, gdzie wzlatują
\VUgras{skowronki} \textsuperscript{(41)}, panuje ożywienie. Wszyscy są zajęci żniwami.
Ale podczas gdy biedni nadal używają dawnych narzędzi, \VUgras{kos, sierpów} i
\VUgras{cepów} \textsuperscript{(40)}, bogaci właściciele każą kosić swoje zboże albo
swoje \VUgras{żyto} \textsuperscript{(43)} \VUgras{żniwiarką} \textsuperscript{(42)}.
Potem, nie musząc znosić snopów do stodoły, robi się wielkie \VUgras{sterty snopów}
\textsuperscript{(44)}.

%%K t08-c1-11-1 p
11. --- Wtedy przywozi się \VUgras{młockarnię} \textsuperscript{(47)}, którą
\VUgras{lokomobila} \textsuperscript{(45)} porusza \VUgras{pasem transmisyjnym}
\textsuperscript{(46)}, i tak ziarno oddziela się od \VUgras{słomy}, nie opuszczając
pola, na którym zostało zasiane.

%%K t08-tit-4 sub
\VUsaut{5.0mm}
\VUpk{10.2pt}{40}{Burza.}
\VUsaut{3.4mm}

%%K t08-c1-12-1 p
12. --- Często \VUgras{nawałnice} \textsuperscript{(53)} zdarzają się w porze żniw.
Najpierw z daleka słychać pomruki \VUgras{grzmotu}; \VUgras{wichura z zachodu}
\textsuperscript{(54)} zaczyna dąć i, dysząc, zgina najgrubsze drzewa \VUgras{lasku}
\textsuperscript{(49)}. Czarne \VUgras{chmurzyska} \textsuperscript{(56)} szturmują
niebo; przecinają je oślepiające zygzaki \VUgras{błyskawicy} \textsuperscript{(55)}.
\VUgras{Deszcz}, a czasem \VUgras{grad} zaczynają padać, przygniatając plony gotowe do
skoszenia.

%%K t08-c1-13-1 p
13. --- Często piorun podpala budowlę. Tak w zeszłym roku dach \VUgras{wiatraka}
\textsuperscript{(50)} obrócił się w popiół, a dwa jego \VUgras{skrzydła}
\textsuperscript{(51)} zostały strzaskane.

%%K t08-c1-14-1 p
14. --- Po kilku godzinach spokój się odradza. Lśniąca \VUgras{tęcza}
\textsuperscript{(52)} pojawia się na widnokręgu, a przyroda podejmuje swoje życie
przerwane na kilka chwil. Wieśniacy, którzy schronili się w \VUgras{szałasach}
\textsuperscript{(48)}, znów zabierają się do pracy i dzielnie starają się naprawić
szkody, które właśnie ponieśli.

%%K t08-tit-5 sub
\VUsaut{6.0mm}
\VUcentre{13.2pt}{90}{\textit{Scena druga.}}
\VUsaut{3.6mm}

%%K t08-tit-6 sub
\VUpk{10.2pt}{40}{Polowanie. --- Winobranie.}
\VUsaut{1.4mm}
\VUpk{10.2pt}{40}{Rybołówstwo.}
\VUsaut{4.0mm}

%%K t08-c2-01-1 p
1. --- Około końca \VUgras{sierpnia} albo w połowie \VUgras{września}, zależnie od
okolic, przypada otwarcie polowania. Ledwie pierwsze promienie słońca wywabiły
\VUgras{jaszczurki} \textsuperscript{(2)} z ich nory, kiedy \VUgras{myśliwy}
\textsuperscript{(5)} wyrusza w drogę ze swoimi \VUgras{psami} \textsuperscript{(4)}.

%%K t08-c2-02-1 p
2. --- Włożył swój mocny \VUgras{strój myśliwski} \textsuperscript{(9)} z grubego sukna,
obuł się w skórzane \VUgras{getry}, w których zdoła chodzić po wilgotnej trawie, gdzie
kryją się \VUgras{ropuchy} \textsuperscript{(1)}; wziął swoją \VUgras{strzelbę}
\textsuperscript{(6)} i swoją \VUgras{torbę myśliwską} \textsuperscript{(10)}, i oto
rusza.

%%K t08-c2-03-1 p
3. --- Zapał pościgu prowadzi go w środek winnicy, gdzie winobranie jest w pełnym
ruchu. Dzieci winiarza, które zwykle pilnują kóz na drodze, dziś pozwalają im paść się
na chybił trafił, a przywiązawszy do palika starego \VUgras{kozła} \textsuperscript{(3)},
który nie omieszkałby skorzystać ze swojej wolności, żeby uciec, one z kolei przypisują
sobie udział w przyjemności. Jedno chwyta kiść winogron w \VUgras{koszu do zbioru}
\textsuperscript{(12)} i bawi się wyciskaniem \VUgras{soku} \textsuperscript{(14)} do
kubka \textsuperscript{(13)}, żeby zrobić moszcz (wino niesfermentowane). Drugie stroi
sobie głowę \VUgras{wieńcem z liści} \textsuperscript{(15)}, który dopiero co uplotło.

Obok nich bezwstydne \VUgras{wróble} \textsuperscript{(16)} przylatują nawet przed prasę,
żeby kraść winogrona.

%%K t08-c2-04-1 p
4. --- Podczas gdy dzieci się bawią, mężczyźni pracują. \VUgras{Winobrańcy}
\textsuperscript{(18)} znoszą winogrona do piwnicy w \VUgras{koszach na plecach}
\textsuperscript{(17)}; \VUgras{bednarz} \textsuperscript{(19)} naprawia \VUgras{beczki}
\textsuperscript{(20)}, sprawdza, czy \VUgras{klepki} \textsuperscript{(22)} i \VUgras{dna}
\textsuperscript{(21)} są w dobrym stanie, i wymienia pęknięte \VUgras{obręcze}
\textsuperscript{(23)}. Pracował też przy \VUgras{kadziach} \textsuperscript{(25)}, do
których wsypuje się \VUgras{winogrona} \textsuperscript{(26)}, żeby je poddać
fermentacji.

%%K t08-c2-05-1 p
5. --- Kiedy fermentacja się skończy, ściąga się wino, a resztki kładzie się na
\VUgras{prasie} \textsuperscript{(28)} i stopniowo dociskając wielką \VUgras{śrubą}
\textsuperscript{(29)}, wyciska się sok, który spływa do \VUgras{kadzi}
\textsuperscript{(32)}, i jeszcze otrzymuje się \VUgras{wino} \textsuperscript{(31)}.

%%K t08-tit-8 sub
\VUsaut{6.0mm}
\VUpk{10.2pt}{40}{Piwo i jabłecznik.}
\VUsaut{3.4mm}

%%K t08-c2-06-1 p
6. --- Po murze \VUgras{piwnicy} \textsuperscript{(27)} pnie się gałąź \VUgras{chmielu}
\textsuperscript{(30)}. Tam jest to tylko roślina ozdobna, ale w niektórych krajach,
gdzie winorośl jest bardzo rzadka, chmiel uprawia się, żeby wyrabiać \VUgras{piwo}.

%%K t08-c2-07-1 p
7. --- Mała \VUgras{jabłoń} \textsuperscript{(33)} jest obładowana jabłkami, które, kiedy
dojrzeją, dadzą wyborny \VUgras{jabłecznik}. W swojej \VUgras{klatce}
\textsuperscript{(34)} \VUgras{kanarek} \textsuperscript{(35)} i \VUgras{szczygieł}
\textsuperscript{(36)} patrzą zazdrosnym okiem na wróble, które swobodnie dokazują.

%%K t08-c2-08-1 p
8. --- Jak okiem sięgnąć, na stoku wzgórz ustawione są \VUgras{tyczki winne}
\textsuperscript{(40)}, które podtrzymują przepiękne \VUgras{krzewy winne}
\textsuperscript{(39)}, obwieszone ciężkimi \VUgras{kiściami}.

%%K t08-c2-08-2 p
Przez cały rok \VUgras{winiarz} \textsuperscript{(42)} pracował, żeby doglądać swojej
\VUgras{winnicy} \textsuperscript{(38)}, a teraz jest wynagrodzony za swój trud pięknym
zbiorem. \VUgras{Winobranki} \textsuperscript{(41)} napełniają kadzie postawione na wozie
ciągnionym przez \VUgras{muły} \textsuperscript{(43)}, i wszyscy są radośni.

%%K t08-tit-9 sub
\VUsaut{5.0mm}
\VUpk{10.2pt}{40}{Rybołówstwo.}
\VUsaut{3.4mm}

%%K t08-c2-09-1 p
9. --- Tak samo jest z \VUgras{wędkarzami} \textsuperscript{(44)}, którzy od rana
przyszli usadowić się na brzegu wody ze swoimi \VUgras{wędkami}
\textsuperscript{(45)}.

%%K t08-c2-09-2 p
\VUgras{Ryba} \textsuperscript{(47)} chwyta \VUgras{haczyk}, gdy tylko zanurzą swoją
\VUgras{żyłkę} \textsuperscript{(46)} w wodzie, i ledwie mają czas odnowić
\VUgras{przynętę}, kiedy nowa ofiara trzepocze na końcu żyłki.

%%K t08-c2-09-3 p
Jednak \VUgras{rybak z siecią} \textsuperscript{(48)}, który ze swojej \VUgras{łodzi}
\textsuperscript{(49)} rzuca \VUgras{sieć} w środek \VUgras{rzeki} \textsuperscript{(50)},
łowi więcej ryb.

%%K t08-c2-10-1 p
10. --- Jesień jest porą dobrych przechadzek: pieszo, \VUgras{konno}
\textsuperscript{(52)}, \VUgras{na rowerze} \textsuperscript{(53)},
\VUgras{samochodem} \textsuperscript{(54)}. \VUgras{Drogi} \textsuperscript{(51)} są
czyste i korzysta się z ostatnich pięknych dni, bo oto już \VUgras{jaskółki}
\textsuperscript{(56)} zbierają się na dachach, żeby odlecieć z całych skrzydeł ku
cieplejszym krajom. Listowie starego \VUgras{orzecha} \textsuperscript{(57)} żółknie,
\VUgras{orzechy} spadają, a wysokie \VUgras{drzewa} ustawione w grupie jak
\VUgras{bukiet} \textsuperscript{(58)} na \VUgras{wzniesieniu} \textsuperscript{(59)}
wkrótce stracą liście.

%%K t08-c2-11-1 p
11. --- \VUgras{Słońce} \textsuperscript{(60)} wstaje później, dni się skracają, dość
szczypiące zimno zaczyna być odczuwane \VUgras{rankiem}, i oto już stada \VUgras{słonek}
\textsuperscript{(61)} opuszczają nasze niegościnne strefy.
\VUsaut{6.0mm}
\VUfilet{22mm}
